\section{Conclusion}
\label{sec:conclusion}

The contribution is a recovery-aware physical attribution construction for sparse PMUs: a restored operator leaves a dynamic state that continues to explain measurements and must compete with later interventions. The implemented load/line bank predicts intervention-conditioned DAE histories and mixes fixed network hypotheses. $M$ records the source masses; $\Delta M$ separates within-model evolution from model reweighting. Physical templates and unknown-onset scans have close precedents, and this study does not establish superiority over unexecuted trajectory-based methods.

The completed development evidence shows both utility and a limit. Physical state recovery can be accurate while weak event sources remain ambiguous. In causal H04 development, the five-model reference obtains \ContMixTimely{}/\ContMixTotal{} timely joint decisions with \ContMixExtras{} extras; nominal improves timeliness at a higher alarm burden. The family and temporal breakdowns explain that result, including the BUS31 delay and the negative BUS16 admission experiment. Final performance claims depend on the frozen confirmatory results. No broader all-family, unseen-source or field claim follows from these tests.

\section*{Supplementary Material}
The supplement preserves detailed proofs, electrical and instrumental conventions, the original batch estimator and joint-estimator proposals, expanded component studies, and complete evidence contracts. The pre-revision manuscript is archived in the repository. Empirical data and inference baselines are unchanged.
